\setcounterpageref{ex}{0}
\TaoPhieuKiemTra{19}

\setcounter{ex}{0}
\setcounter{bt}{0}

\textbf{\large I. PHẦN TRẮC NGHIỆM (3,0 điểm)}

{\itshape Khoanh tròn vào chữ cái đặt trước câu trả lời đúng. Mỗi câu đúng được 0,5 điểm.}

\begin{ex} % Câu 1
Trên tia số, điểm nằm giữa số $28$ và số $30$ biểu diễn số
\choice
{\True $29$}
{$27$}
{$31$}
{$30$}
\loigiai{
Số nằm giữa $28$ và $30$ là $29$.\\
\textbf{Đáp án đúng: A}
}
\end{ex}

\begin{ex} % Câu 2
Cho phép tính $62 + 25 = 87$. Số $87$ được gọi là
\choice
{Số hạng}
{Số trừ}
{\True Tổng}
{Hiệu}
\loigiai{
Số $87$ là kết quả phép cộng nên gọi là Tổng.\\
\textbf{Đáp án đúng: C}
}
\end{ex}

\begin{ex} % Câu 3
Tính: $14\text{ kg} - 8\text{ kg} + 5\text{ kg} = \dots$
\choice
{$10\text{ kg}$}
{\True $11\text{ kg}$}
{$12\text{ kg}$}
{$1\text{ kg}$}
\loigiai{
$14 - 8 = 6\text{ kg}$; $6 + 5 = 11\text{ kg}$.\\
\textbf{Đáp án đúng: B}
}
\end{ex}

\begin{ex} % Câu 4
Hình tròn có bao nhiêu cạnh thẳng?
\choice
{$1$ cạnh}
{$2$ cạnh}
{$4$ cạnh}
{\True $0$ cạnh}
\loigiai{
Hình tròn được bao quanh bởi đường cong khép kín, không có cạnh thẳng.\\
\textbf{Đáp án đúng: D}
}
\end{ex}

\begin{ex} % Câu 5
So sánh: $12 - 5 \dots 11 - 3$.
\choice
{\True $<$}
{$>$}
{$=$}
{Không so sánh được}
\loigiai{
$12 - 5 = 7$; $11 - 3 = 8$. Vì $7 < 8$ nên điền dấu $<$.\\
\textbf{Đáp án đúng: A}
}
\end{ex}

\begin{ex} % Câu 6
Tính: $9 + 9 - 10 = \dots$
\choice
{$7$}
{\True $8$}
{$9$}
{$10$}
\loigiai{
$9 + 9 = 18$; $18 - 10 = 8$.\\
\textbf{Đáp án đúng: B}
}
\end{ex}

\vspace{4mm}
\textbf{\large II. PHẦN TỰ LUẬN (7,0 điểm)}

{\itshape Thí sinh trình bày chi tiết bài làm vào giấy kiểm tra.}

\begin{bt}[1,0 điểm]
Tính nhẩm:
\begin{multicols}{2}
\begin{enumerate}[a)]
  \item $9 + 8$
  \item $7 + 9$
  \item $16 - 8$
  \item $18 - 9$
\end{enumerate}
\end{multicols}
\loigiai{
\begin{itemize}
  \item $9 + 8 = 17$
  \item $7 + 9 = 16$
  \item $16 - 8 = 8$
  \item $18 - 9 = 9$
\end{itemize}
}
\end{bt}

\begin{bt}[1,5 điểm]
Đặt tính rồi tính:
\begin{multicols}{3}
\begin{enumerate}[a)]
  \item $31 + 39$
  \item $93 - 38$
  \item $54 + 38$
\end{enumerate}
\end{multicols}
\loigiai{
\begin{itemize}
  \item $31 + 39 = 70$
  \item $93 - 38 = 55$
  \item $54 + 38 = 92$
\end{itemize}
}
\end{bt}

\begin{bt}[1,0 điểm]
Tính giá trị biểu thức kèm đơn vị đo:
\begin{enumerate}[a)]
  \item $35\text{ kg} + 10\text{ kg} - 10\text{ kg}$
  \item $44\text{ }l - 15\text{ }l + 9\text{ }l$
\end{enumerate}
\loigiai{
\begin{enumerate}[a)]
  \item $35\text{ kg} + 10\text{ kg} - 10\text{ kg} = 45\text{ kg} - 10\text{ kg} = 35\text{ kg}$.
  \item $44\text{ }l - 15\text{ }l + 9\text{ }l = 29\text{ }l + 9\text{ }l = 38\text{ }l$.
\end{enumerate}
}
\end{bt}

\begin{bt}[1,5 điểm]
Bạn Tùng có $34$ con tem. Bạn Hoa có nhiều hơn bạn Tùng $5$ con tem. Hỏi bạn Hoa có bao nhiêu con tem?
\loigiai{
Số con tem bạn Hoa có là:
\[ 34 + 5 = 39 \text{ (con tem)} \]
\textbf{Đáp số:} $39$ con tem.
}
\end{bt}

\begin{bt}[1,0 điểm]
Một thùng chứa đầy $54\text{ }\text{kg}$ đường kính. Buổi sáng cửa hàng đã bán ra $15\text{ }\text{kg}$. Hỏi trong thùng còn lại bao nhiêu ki-lô-gam đường kính?
\loigiai{
Trong thùng còn lại số ki-lô-gam đường kính là:
\[ 54 - 15 = 39\text{ }\text{kg} \]
\textbf{Đáp số:} $39\text{ }\text{kg}$ đường kính.
}
\end{bt}

\begin{bt}[1,0 điểm]
Tìm số bé nhất có hai chữ số mà tổng hai chữ số của nó bằng $9$.
\loigiai{
Để số bé nhất thì chữ số hàng chục bé nhất, chọn chữ số hàng chục là $1$.\\
Chữ số hàng đơn vị: $9 - 1 = 8$.\\
Vậy số đó là $18$.\\
\textbf{Đáp số:} $18$.
}
\end{bt}
